\documentclass[12pt,a4paper]{report}

\usepackage[T1]{fontenc}
\usepackage[utf8]{inputenc}
\usepackage[french]{babel}
\usepackage{geometry}
\usepackage{setspace}
\usepackage{lmodern}
\usepackage{tikz}
\usepackage{xcolor}

\geometry{
    top=2.5cm,
    bottom=2.5cm,
    left=2.5cm,
    right=2.5cm
}

\setstretch{1.15}

\begin{document}

%==================================
% TITRE PERSONNALISÉ
%==================================

\vspace*{0.5cm}

\noindent
\begin{tikzpicture}

% Rectangle noir
\fill[black] (0,0) rectangle (0.9,1.8);

% Ligne horizontale
\draw[line width=0.5pt] (1.2,0.9) -- (14.8,0.9);

% Titre
\node[
anchor=west,
font=\fontsize{24}{28}\bfseries
] at (1.2,0.25)
{TABLE DES MATIÈRES};

\end{tikzpicture}

\vspace{1.2cm}

\tableofcontents
\clearpage

%=================================================
% INTRODUCTION GENERALE
%=================================================

\chapter*{Introduction générale}
\addcontentsline{toc}{chapter}{Introduction générale}

Texte introduction...

%=================================================
% CHAPITRE 1
%=================================================

\chapter{Cadre général du projet}

\section{Introduction}

\section{Présentation de l'organisme d'accueil}

\section{Présentation du projet}
\subsection{Contexte du projet}
\subsection{Problématique}
\subsection{Étude critique de l'existant}
\subsection{Solution proposée}

\section{Méthodologie de travail}
\subsection{Méthodologie Agile Scrum}
\subsection{Planification globale du projet}

\section{Environnement technique}
\subsection{Technologies backend}
\subsection{Technologies frontend}
\subsection{Base de données et accès aux données}
\subsection{Outils de développement et de test}

\section{Conclusion}

%=================================================
% CHAPITRE 2
%=================================================

\chapter{Analyse et spécification des besoins}

\section{Introduction}

\section{Identification des acteurs}
\subsection{Super administrateur}
\subsection{Administrateur}
\subsection{Utilisateur final}

\section{Clarification du vocabulaire métier}
\subsection{Famille de documents}
\subsection{Modèle de document ou template}
\subsection{Charte graphique}
\subsection{Configuration de vue de table}
\subsection{Instance de document}

\section{Analyse des besoins fonctionnels}
\subsection{Gestion des utilisateurs}
\subsection{Gestion des rôles et permissions}
\subsection{Gestion des familles de documents}
\subsection{Gestion des modèles de documents}
\subsection{Gestion des chartes graphiques}
\subsection{Gestion des configurations de vues de tables}
\subsection{Édition et génération des documents}
\subsection{Gestion des bénéficiaires et des données métier}

\section{Analyse des besoins non fonctionnels}
\subsection{Sécurité}
\subsection{Performance}
\subsection{Ergonomie et expérience utilisateur}
\subsection{Maintenabilité}

\section{Backlog du produit}
\subsection{User stories principales}
\subsection{Priorisation des besoins}
\subsection{Critères d'acceptation}

\section{Planification des sprints}
\subsection{Découpage temporel}
\subsection{Objectifs par sprint}
\subsection{Livrables attendus}

\section{Diagrammes UML globaux}
\subsection{Diagramme de cas d'utilisation général}
\subsection{Diagramme de classes global}
\subsection{Diagramme de séquence global}

\section{Conclusion}

%=================================================
% CHAPITRE 3
%=================================================

\chapter{Sprint 1 : Authentification et socle applicatif}

\section{Introduction}
\section{Objectifs du sprint}
\section{Backlog du sprint}
\subsection{Authentification JWT}
\subsection{Gestion du profil utilisateur}
\subsection{Configuration du backend}
\subsection{Connexion à la base de données}
\section{Analyse et conception}
\subsection{Diagramme de cas d'utilisation}
\subsection{Diagramme de classes}
\subsection{Diagrammes de séquence}
\section{Réalisation}
\subsection{Mise en place du backend ASP.NET Core}
\subsection{Implémentation de l'authentification}
\subsection{Configuration CORS, Swagger et middlewares}
\subsection{Configuration de la base SQL Server}
\section{Tests et validation}
\section{Conclusion}

%=================================================
% CHAPITRE 4
%=================================================

\chapter{Sprint 2 : Gestion des rôles et gouvernance d'accès}

\section{Introduction}
\section{Objectifs du sprint}
\section{Backlog du sprint}
\subsection{Définition des rôles}
\subsection{Protection des routes et des endpoints}
\subsection{Gestion des utilisateurs et profils}
\section{Analyse et conception}
\subsection{Diagramme de cas d'utilisation}
\subsection{Diagramme de classes}
\subsection{Diagrammes de séquence}
\section{Réalisation}
\subsection{Implémentation du système RBAC}
\subsection{Mise en place des autorisations backend}
\subsection{Gestion des guards côté frontend}
\subsection{Sécurisation des accès par profil}
\section{Tests et validation}
\section{Conclusion}

%=================================================
% CHAPITRE 5
%=================================================

\chapter{Sprint 3 : Module coeur d'édition documentaire}

\section{Introduction}
\section{Objectifs du sprint}
\section{Backlog du sprint}
\subsection{Intégration de l'éditeur riche}
\subsection{Création et modification du contenu documentaire}
\subsection{Sauvegarde et chargement des contenus}
\subsection{Prévisualisation du document}
\section{Analyse et conception}
\subsection{Diagramme de cas d'utilisation}
\subsection{Diagramme de classes}
\subsection{Diagrammes de séquence}
\section{Réalisation}
\subsection{Intégration de l'éditeur Tiptap}
\subsection{Gestion du contenu des modèles de documents}
\subsection{Application des en-têtes, corps et pieds de page}
\subsection{Prévisualisation et rendu des documents}
\section{Tests et validation}
\section{Conclusion}

%=================================================
% CHAPITRE 6
%=================================================

\chapter{Sprint 4 : Administration des modules métier}

\section{Introduction}
\section{Objectifs du sprint}
\section{Backlog du sprint}
\subsection{Gestion des familles de documents}
\subsection{Gestion des modèles de documents}
\subsection{Gestion des chartes graphiques}
\subsection{Gestion des configurations de vues de tables}
\subsection{Gestion des bénéficiaires et filtres métier}
\section{Analyse et conception}
\subsection{Diagramme de cas d'utilisation}
\subsection{Diagramme de classes}
\subsection{Diagrammes de séquence}
\section{Réalisation}
\subsection{CRUD des familles de documents}
\subsection{CRUD des modèles de documents}
\subsection{CRUD des chartes graphiques}
\subsection{CRUD des configurations de vues de tables}
\subsection{Association famille, modèle et charte graphique}
\subsection{Configuration des données métier et des filtres}
\section{Tests et validation}
\section{Conclusion}

%=================================================
% CHAPITRE 7
%=================================================

\chapter{Sprint 5 : Intégration, qualité et livraison}

\section{Introduction}
\section{Objectifs du sprint}
\section{Backlog du sprint}
\subsection{Intégration frontend/backend}
\subsection{Tests fonctionnels et tests d'intégration}
\subsection{Optimisation des performances}
\subsection{Préparation de la livraison}
\section{Réalisation}
\subsection{Validation des parcours utilisateurs}
\subsection{Gestion des erreurs}
\subsection{Optimisation et renforcement de la sécurité}
\subsection{Documentation technique et utilisateur}
\subsection{Préparation de la démonstration PFE}
\section{Tests et validation}
\section{Conclusion}

%=================================================
% CONCLUSION
%=================================================

\chapter*{Conclusion générale}
\addcontentsline{toc}{chapter}{Conclusion générale}

Texte conclusion...

%=================================================
% WEBOGRAPHIE
%=================================================

\chapter*{Webographie}
\addcontentsline{toc}{chapter}{Webographie}

%=================================================
% BIBLIOGRAPHIE
%=================================================

\chapter*{Bibliographie}
\addcontentsline{toc}{chapter}{Bibliographie}

%=================================================
% ANNEXES
%=================================================

\chapter*{Annexes}
\addcontentsline{toc}{chapter}{Annexes}

\end{document}
